\section{Introduction}
Ophthalmic assessment requires interpreting spatial abnormalities together with the clinical circumstances in which an image was acquired. In ocular toxoplasmosis, lesion assessment is part of a broader diagnostic and follow-up process; a survey of international practice documents the roles of clinical examination, investigations, treatment decisions, and monitoring~\cite{yogeswaran2023}. An automated segmentation system can provide a reproducible image-space representation of annotated abnormalities. However, a segmented region alone does not establish infectious activity, identify the biological origin of disease, or predict the time to visual deterioration. These distinctions determine the scope of the present study.

The Ocular Toxoplasmosis Fundus Images Dataset (OTFID) provides fundus photographs, lesion masks, and image-level labels~\cite{parra2021}. It offers an accessible starting point for algorithm development, but its released annotation table does not provide the verified patient linkage, temporally aligned electronic medical records (EMRs), or longitudinal visual outcomes required by the full OcuWeave research agenda. Our release-specific audit also identifies nonhealthy photographs without eligible matching masks. Treating those photographs as healthy would introduce incorrect supervision. We therefore formulate the current task as binary segmentation on explicitly eligible image--mask pairs and retain the remaining photographs in the source archive.

Three methodological issues motivate the experimental design. First, resizing changes both image detail and the sampled reference boundary. In particular, survival of a nonempty union mask does not establish survival of every connected foreground component. Second, overlap measurements can obscure different types of error: false-positive area on a healthy image and undersegmentation of a small annotated region have different interpretations. Problem-aware validation therefore requires explicit metric definitions and complementary measurements~\cite{maierhein2024}. Third, unknown patient identities and visually similar images constrain the independence of development and evaluation sets. Technical similarity grouping reduces one identifiable source of overlap, but cannot substitute for verified patient, eye, and visit identifiers.

Modern convolutional designs demonstrate the continued relevance of spatial and channel mixing~\cite{liu2022}, while query-based mask prediction provides an alternative representation for segmentation~\cite{cheng2022}. Building on these architectural ideas, we implement OcuWeave--LQNet, a compact Local-detail Query Network. Its base model projects multiscale convolutional features onto one pixel grid and combines learned query masks into a binary logit map. A shallow image-conditioned branch provides an additive local correction. The implementation uses randomly initialized PyTorch modules; it does not use a U-Net architecture, a pretrained checkpoint, or a reproduced Mask2Former decoder. Architectural provenance is distinguished from a claim of algorithmic novelty.

The immediate study asks two bounded questions: does increasing the input size from 192 to 384 pixels improve original-grid validation segmentation under a fixed update budget, and does local refinement provide an additional improvement at 384 pixels? We specify paired seeds, immutable data and configuration receipts, and separate preprocessing diagnostics from model evaluation. The contributions at this stage are an inspectable implementation, an auditable development protocol, and a prespecified comparison design. Performance gains remain hypotheses until the corresponding runs are completed.

The broader vision--language objective is deliberately deferred. Harvard-FairVLMed supplies a relevant image--clinical-note resource for a separate glaucoma investigation~\cite{luo2024}; it does not supply EMRs for OTFID patients. No cross-dataset patient pairing is performed. Future multimodal work will require verified within-cohort pairing, temporal availability checks, and target-leakage review before a learning objective or clinical claim can be justified.
